%% kapa: exam=1 round=32 year=2021 period=1 form=A subject=real_estate subject_ko=부동산학원론 seq=36 number=116
%% source: Q-Net 감정평가사 기출문제 (cst00309) artlSeq=5208397 file=extracted/1008/5208397/A형/2021년도 제32회 감정평가사 1차 1교시 A형.pdf
%% conversion: kordoc 4.24.51 -> tools/gen1.py (auto); review=proofread-sonnet; flags=none
\begin{question}{116}
감정평가에 관한 규칙상 현황기준 원칙에 관한 내용으로 옳지 않은 것은? (단, 감정평가조건이란 기준시점의 가치형성요인 등을 실제와 다르게 가정하거나 특수한 경우로 한정하는 조건을 말함)
\begin{choices}
  \item 감정평가업자는 감정평가조건의 합리성, 적법성이 결여되거나 사실상 실현 불가능하다고 판단할 때에는 의뢰를 거부하거나 수임을 철회할 수 있다.
  \item 현황기준 원칙에도 불구하고 법령에 다른 규정이 있는 경우에는 감정평가조건을 붙여 감정평가할 수 있다.
  \item 현황기준 원칙에도 불구하고 대상물건의 특성에 비추어 사회통념상 필요하다고 인정되는 경우에는 감정평가조건을 붙여 감정평가할 수 있다.
  \item 감정평가의 목적에 비추어 사회통념상 필요하다고 인정되어 감정평가조건을 붙여 감정평가하는 경우에는 감정평가조건의 합리성, 적법성 및 실현가능성의 검토를 생략할 수 있다.
  \item 현황기준 원칙에도 불구하고 감정평가 의뢰인이 요청하는 경우에는 감정평가조건을 붙여 감정평가할 수 있다.
\end{choices}
\end{question}
\answer{4}
